\documentclass[12pt,a4paper]{article}

% Encoding and fonts
\usepackage[utf8]{inputenc}
\usepackage[T1]{fontenc}
\usepackage{lmodern}

% Mathematics
\usepackage{amsmath}
\usepackage{amssymb}
\usepackage{amsfonts}
\usepackage{amsthm}
\usepackage{mathtools}

% Layout
\usepackage[margin=1in]{geometry}
\usepackage{setspace}
\doublespacing
\emergencystretch=2em

% References and links
\usepackage[hidelinks]{hyperref}

% Figures
\usepackage{graphicx}
\graphicspath{{figures/}}
\usepackage{float}
\usepackage[hypcap=true]{caption}

% Theorem environments
\newtheorem{theorem}{Theorem}[section]
\newtheorem{proposition}[theorem]{Proposition}
\newtheorem{lemma}[theorem]{Lemma}
\newtheorem{corollary}[theorem]{Corollary}
\theoremstyle{definition}
\newtheorem{definition}[theorem]{Definition}
\theoremstyle{remark}
\newtheorem{remark}[theorem]{Remark}
\newtheorem{conjecture}[theorem]{Conjecture}

% Title information
\title{Decoherence as Dimensional Collapse: The Geometry of the Quantum-Classical Transition}
\author{Ian Todd\\
Sydney Medical School\\
University of Sydney\\
Sydney, NSW, Australia\\
\texttt{itod2305@uni.sydney.edu.au}}
\date{}

\begin{document}

\maketitle

\begin{abstract}
The quantum-to-classical transition via decoherence is well established, but its information-geometric structure has not been quantified in the explicit Bures/Fisher--Rao framework. We show that decoherence is a dimensional collapse in quantum state space: the dephasing channel maps a $(d^2-1)$-dimensional quantum manifold onto the $(d-1)$-dimensional classical probability simplex, annihilating the $d^2-d$ coherence directions. At full-rank diagonal states, the classical and quantum tangent subspaces are Bures-orthogonal, so the differential of the dephasing map is the metric-orthogonal projector onto the classical subspace. As our main result, we compute the exact Bures principal angle $\theta_B$ between these subspaces at every full-rank qubit state $(|\mathbf{r}| < 1)$:
\[
\cos^2\theta_B = \frac{r_z^2\, r_\perp^2}{(1-r_\perp^2)(1-r_z^2)},
\]
where $r_z$ is the population imbalance and $r_\perp$ the coherence magnitude. The angle is $\pi/2$ (perfect orthogonality) when the state is diagonal or has equal populations; the split is oblique exactly when both coherence and population imbalance are present, and dephasing monotonically restores orthogonality. The separation angle has an operational meaning: $\sin^2\theta_B$ is the fraction of SLD Fisher information for the population parameter that survives Schur-complement elimination of coherence nuisance parameters. For $d \geq 3$, we prove that Bures-orthogonality fails whenever a nonzero coherence has a decoupled spectator index, and give a perturbative cross-Gram formula quantifying the departure near the classical simplex.  We conjecture that for $d \geq 3$, only diagonal states achieve orthogonality, supported by numerical evidence for $d = 3, 4, 5$. When the Hamiltonian is diagonal in the dephasing basis, coherences store non-equilibrium free energy $k_BT \cdot C_r$, which decoherence irreversibly destroys. These results use standard quantum mechanics alone and connect Zurek's decoherence program to information geometry.
\end{abstract}

\noindent\textbf{Keywords:} decoherence, quantum Fisher information, information geometry, dimensional collapse, quantum-classical transition

\begin{remark}[On ``dimensionality'']
\label{rem:dimensionality}
Throughout this paper we use \emph{dimensionality} in the \textbf{information-geometric sense}: the effective degrees of freedom of distinguishable dynamics, quantified by the rank or participation ratio of the quantum Fisher information matrix---\textbf{not} the Hilbert space dimension $d$ of the system. Decoherence is dimensional collapse: the reduction of information-geometric dimensionality from $d^2 - 1$ (full quantum) to $d - 1$ (classical).
\end{remark}

%% ============================================================
\section{Introduction}
\label{sec:intro}

The decoherence program, developed by Zeh \cite{zeh1970interpretation} and Zurek \cite{zurek1981pointer,zurek2003decoherence}, explains how classical definiteness emerges from quantum superpositions without requiring interpretive commitments beyond standard quantum mechanics \cite{schlosshauer2005decoherence}.  When a quantum system interacts with its environment, entanglement causes the system's coherences---the off-diagonal elements of its density matrix in the pointer basis---to decay on extremely short timescales (${\sim}\,10^{-31}$\,s for a dust grain in a photon bath \cite{zurek2003decoherence}).  Zurek's quantum Darwinism \cite{zurek2025decoherence,riedel2010zurek} extends this: pointer states that survive decoherence are replicated across independent environmental fragments, creating the redundant classical records that underwrite objectivity.

Decoherence theory explains \emph{that} the quantum-to-classical transition occurs, but not \emph{how much structure it destroys}.  How many degrees of freedom does a quantum system lose when it decoheres?  What is the precise dimensional structure of the map from quantum to classical?  How cleanly does the surviving classical information separate from the destroyed quantum information?  These questions have natural answers in information geometry---the study of statistical manifolds equipped with the Fisher information metric \cite{rao1945information,chentsov1982statistical,amari1985differential}---but this connection has not been made explicit for decoherence in the Bures/Fisher--Rao framework.  The quantum Fisher information matrix (QFIM) \cite{liu2020quantum,braunstein1994statistical} measures the dimensionality of distinguishable quantum states, and the Bures metric it induces \cite{bures1969extension,uhlmann1976transition,petz1996monotone} is the unique minimal monotone Riemannian metric on quantum state space.  These tools are well established in quantum estimation theory, but they have not been applied systematically to decoherence as a question about dimensional structure.

In this paper, we show that decoherence is a \emph{dimensional collapse} in the information-geometric sense. A $d$-dimensional quantum system lives on a $(d^2-1)$-dimensional state manifold $\mathcal{M}_Q$ equipped with the Bures metric. Decoherence projects this manifold onto the $(d-1)$-dimensional classical probability simplex $\Delta_{d-1}$ equipped with the Fisher--Rao metric. Specifically:
\begin{enumerate}
\item[\emph{(Background.)}] The decoherence channel reduces the QFIM rank from $d^2-1$ to at most $d-1$ (Theorem~\ref{thm:dim-collapse}), the Bures metric restricts to the Fisher--Rao metric on $\Delta_{d-1}$ (Theorem~\ref{thm:metric-restriction}), and at full-rank diagonal states the differential of the dephasing map is the Bures-orthogonal projector onto the classical tangent subspace (Theorem~\ref{thm:orthogonal-decomposition}).
\item[\emph{(Main result.)}] We compute the exact Bures principal angle $\theta_B$ between the classical and quantum tangent subspaces at every full-rank qubit state (Theorem~\ref{thm:bures-angle}).  Orthogonality holds when the state is diagonal or has equal populations; the split is oblique exactly when both coherence and population imbalance are present, and dephasing monotonically restores it.  The angle has an operational meaning: $\sin^2\theta_B$ is the fraction of SLD Fisher information for the population parameter that survives elimination of the coherence nuisance parameters (Corollary~\ref{cor:estimation}).
\item[\emph{(Generalization.)}] For $d \geq 3$, whenever a nonzero off-diagonal entry has a decoupled spectator index, Bures-orthogonality fails (Theorem~\ref{thm:spectator-orthogonality}); a perturbative cross-Gram formula (Theorem~\ref{thm:qudit-angles}) quantifies the departure near the simplex.  We conjecture that for $d \geq 3$, only diagonal states achieve orthogonality (Conjecture~\ref{conj:global-orthogonality}), supported by a numerical survey for $d = 3, 4, 5$.
\item[\emph{(Thermodynamics.)}] When $H$ is diagonal in the dephasing basis, coherences store non-equilibrium free energy $k_B T \cdot C_r$, which decoherence irreversibly destroys (Proposition~\ref{prop:free-energy-coherence}).
\end{enumerate}
These results use standard quantum mechanics---no new assumptions, no modifications to the Schr\"odinger equation, no interpretive commitments beyond the Born rule, and no speculative physics are required.

The Born rule is assumed throughout; we do not derive it or explain why a particular outcome is selected.  Our results are compatible with any interpretation that accepts unitary evolution and environmental interaction as the source of classicality.  What we add to the decoherence program is not a new mechanism but a precise geometric characterization of what that mechanism accomplishes.

\paragraph{What is new.}  The dimensional collapse and metric restriction (Theorems~\ref{thm:dim-collapse}--\ref{thm:metric-restriction}) follow from standard properties of quantum channels, and the free energy identity (Proposition~\ref{prop:free-energy-coherence}) rearranges known results \cite{lostaglio2015description,korzekwa2016extraction,streltsov2017colloquium}.  The new contributions are:
\begin{itemize}
\item The orthogonal decomposition (Theorem~\ref{thm:orthogonal-decomposition}), establishing the geometric base case: at diagonal states, the dephasing differential is the Bures-orthogonal projector onto the classical tangent subspace.
\item The exact Bures separation angle (Theorem~\ref{thm:bures-angle}), our main result: a closed-form formula quantifying the departure from orthogonality at every full-rank qubit state, with an operational interpretation via the SLD Schur complement (Corollary~\ref{cor:estimation}).
\item The spectator orthogonality criterion (Theorem~\ref{thm:spectator-orthogonality}), proving that for $d \geq 3$, Bures-orthogonality fails whenever a nonzero coherence has a decoupled spectator index---covering all states with sparse coherence structure and supporting the conjecture that only diagonal states achieve orthogonality (Conjecture~\ref{conj:global-orthogonality}).
\end{itemize}

The paper is organized as follows. Section~\ref{sec:geometry} establishes the information-geometric framework. Section~\ref{sec:decoherence} contains the main results: the dimensional collapse and orthogonal decomposition theorems, the exact Bures separation angle and its estimation-theoretic corollary, the qudit near-simplex theorem, and the spectator orthogonality criterion for $d \geq 3$.  Section~\ref{sec:implications} derives the free energy decomposition and connects the framework to Zurek's einselection and quantum Darwinism.  Section~\ref{sec:examples} gives worked examples.  Section~\ref{sec:discussion} discusses implications and limitations.

%% ============================================================
\section{The Geometry of Quantum State Space}
\label{sec:geometry}

This section establishes the information-geometric framework that makes ``dimensional collapse'' precise. We define the quantum state manifold, equip it with the quantum Fisher information metric, and introduce a participation-ratio measure of effective dimensionality that will serve as the central quantity throughout the paper.

\begin{definition}[Quantum state manifold]
\label{def:state-manifold}
The space of $d \times d$ density matrices,
\begin{equation}
\label{eq:state-manifold}
\mathcal{M}_Q = \bigl\{\rho \in \mathbb{C}^{d \times d} : \rho \geq 0,\; \mathrm{Tr}\,\rho = 1\bigr\},
\end{equation}
is a real manifold of dimension $d^2 - 1$.
\end{definition}

The dimension count follows from a standard parameterization. A $d \times d$ Hermitian matrix has $d^2$ real parameters: $d$ real diagonal entries and $d(d-1)/2$ complex off-diagonal entries in the upper triangle, each contributing two real parameters (real and imaginary parts), giving $d + 2 \cdot d(d-1)/2 = d^2$ in total. The trace constraint $\mathrm{Tr}\,\rho = 1$ removes one degree of freedom, leaving $d^2 - 1$ independent real parameters. Among these, $d - 1$ are diagonal \emph{populations} (eigenvalue parameters) and $d(d-1)$ are off-diagonal \emph{coherences}. The positivity constraint $\rho \geq 0$ restricts the manifold to a convex body whose boundary consists of states with at least one vanishing eigenvalue; the interior consists of full-rank states.

\begin{definition}[Quantum Fisher information matrix]
\label{def:qfim}
Let $\rho(\theta)$ be a smooth parametric family of density matrices with $\theta \in \mathbb{R}^p$. The \emph{quantum Fisher information matrix} (QFIM) is the $p \times p$ matrix with entries
\begin{equation}
\label{eq:qfim}
[F_Q]_{ij} = \mathrm{Tr}\!\left[\rho\,\frac{L_i L_j + L_j L_i}{2}\right],
\end{equation}
where $L_i$ is the \emph{symmetric logarithmic derivative} (SLD) defined implicitly by
\begin{equation}
\label{eq:sld}
\frac{\partial \rho}{\partial \theta^i} = \frac{\rho\, L_i + L_i\, \rho}{2}.
\end{equation}
\end{definition}

The SLD equation~\eqref{eq:sld} has a unique solution whenever $\rho$ is full-rank. In the eigenbasis $\rho = \sum_k p_k \lvert k \rangle\!\langle k \rvert$, the matrix elements of $L_i$ are
\begin{equation}
\label{eq:sld-matrix}
\langle m \rvert L_i \lvert n \rangle = \frac{2}{p_m + p_n}\,\langle m \rvert \partial_i \rho \lvert n \rangle,
\end{equation}
provided $p_m + p_n > 0$. The matrix $F_Q$ is real, symmetric, and positive semidefinite. Its eigenvalues quantify the statistical distinguishability of nearby quantum states along each parameter direction: large eigenvalues correspond to directions in which small parameter changes produce statistically detectable changes in the state, while vanishing eigenvalues correspond to directions along which the state is locally indistinguishable \cite{liu2020quantum,petz2011introduction}.

\begin{proposition}[Bures metric]
\label{prop:bures}
The QFIM defines the Bures--Helstrom metric on~$\mathcal{M}_Q$:
\begin{equation}
\label{eq:bures}
ds^2_{\mathrm{Bures}} = \frac{1}{4}\sum_{i,j} [F_Q]_{ij}\, d\theta^i\, d\theta^j.
\end{equation}
This is the unique minimal monotone Riemannian metric on quantum state space.
\end{proposition}

\begin{proof}[Proof sketch]
Petz's classification theorem \cite{petz1996monotone} establishes a one-to-one correspondence between monotone metrics on $\mathcal{M}_Q$---Riemannian metrics that do not increase distances under any completely positive trace-preserving (CPTP) map---and operator monotone functions $f:(0,\infty) \to (0,\infty)$ satisfying $f(t) = t f(1/t)$. Each such function generates a metric through the formula $\langle A, B \rangle_\rho = \mathrm{Tr}[A\, c_f(L_\rho, R_\rho)^{-1}(B)]$, where $L_\rho$ and $R_\rho$ denote left and right multiplication by $\rho$, and $c_f(x,y) = y f(x/y)$. The SLD metric corresponds to $f(t) = (1+t)/2$, which is the largest operator monotone function in this class. Since a larger $f$ yields a smaller metric tensor, the SLD metric is the \emph{minimal} monotone metric \cite{petz1996monotone}. The connection to the Bures distance \cite{bures1969extension,uhlmann1976transition} was established by Braunstein and Caves \cite{braunstein1994statistical}, building on Wootters' statistical distance \cite{wootters1981statistical}: for infinitesimally separated states, the Bures distance $d_B(\rho, \rho + d\rho)^2 = ds^2_{\mathrm{Bures}}$.
\end{proof}

\begin{proposition}[Data processing inequality]
\label{prop:dpi}
For any CPTP map $\Lambda$,
\begin{equation}
\label{eq:dpi}
F_Q\bigl(\Lambda[\rho(\theta)]\bigr) \leq F_Q\bigl(\rho(\theta)\bigr)
\end{equation}
in the L\"owner (positive semidefinite) order. That is, $F_Q(\rho) - F_Q(\Lambda[\rho])$ is positive semidefinite.
\end{proposition}

\begin{proof}[Proof sketch]
This is a direct consequence of Petz's monotonicity theorem \cite{petz1996monotone}: the Bures metric is monotone under CPTP maps, hence the associated Fisher information cannot increase. An elementary proof using the convexity of the SLD Fisher information under mixing, combined with the Stinespring dilation of $\Lambda$, appears in \cite{liu2020quantum}. For the single-parameter case, a particularly transparent derivation is given by Ferrie \cite{ferrie2014best}.
\end{proof}

The physical content of Proposition~\ref{prop:dpi} is fundamental to what follows: quantum channels cannot increase the information content of a state. Every interaction with an environment---every physical process described by a CPTP map---can only reduce or preserve the statistical distinguishability of quantum states. This is the information-geometric expression of the intuition that ``noise destroys information.'' In Section~\ref{sec:decoherence}, we will show that decoherence channels achieve a specific, structured form of this reduction: they annihilate the coherence directions while preserving the population directions.

\begin{definition}[Effective quantum dimensionality]
\label{def:deff}
Given a QFIM $F_Q$ with eigenvalues $\lambda_1 \geq \lambda_2 \geq \cdots \geq \lambda_p \geq 0$, the \emph{effective quantum dimensionality} is the participation ratio
\begin{equation}
\label{eq:deff}
D_{\mathrm{eff}}^Q(\rho, \theta) = \frac{\bigl(\sum_{i=1}^p \lambda_i\bigr)^2}{\sum_{i=1}^p \lambda_i^2}.
\end{equation}
This ranges from $1$ (a single estimable parameter dominates the Fisher information) to $p$ (all parameters are equally estimable, i.e., all eigenvalues are equal).  The value of $D_{\mathrm{eff}}^Q$ depends on the choice of parameterization.  In general arguments we interpret $D_{\mathrm{eff}}^Q$ relative to a fixed local generator basis: $\rho = I/d + \sum_{a=1}^{d^2-1} \theta^a G_a$, where $\{G_a\}$ is a fixed traceless Hermitian Hilbert--Schmidt-orthonormal basis satisfying $\mathrm{Tr}[G_a G_b] = \delta_{ab}/2$.  In the qubit examples of Section~\ref{sec:qubit-dephasing}, $D_{\mathrm{eff}}^Q$ is computed in the specific angular coordinates $(\theta,\varphi)$ used for that family.
\end{definition}

\begin{remark}[Status and limitations of $D_{\mathrm{eff}}^Q$]
\label{rem:deff-status}
$D_{\mathrm{eff}}^Q$ is a coordinate-dependent diagnostic: the rank of $\mathcal{F}_Q$ is invariant, but the participation ratio depends on the parameterization.  The invariant results of this paper---Theorems~\ref{thm:dim-collapse}--\ref{thm:metric-restriction}, the orthogonal decomposition (Theorem~\ref{thm:orthogonal-decomposition}), and the Bures separation angle (Theorem~\ref{thm:bures-angle})---depend only on rank, kernel structure, and metric properties.  $D_{\mathrm{eff}}^Q$ serves as a continuous illustrative quantity that interpolates between rank bounds in the chosen coordinates; it is used in the worked examples of Section~\ref{sec:examples} but is not needed for the main theorems.  (Note: Loewner monotonicity of the QFIM under CPTP maps does not imply monotonicity of $D_{\mathrm{eff}}^Q$; the invariant content of the data processing inequality is the rank bound.)  $D_{\mathrm{eff}}^Q$ should be distinguished from the Hilbert-space dimensional descriptors of Garc\'{\i}a-P\'{e}rez et al.\ \cite{garcia2023descriptors}, which use eigenvalues of $\rho$ rather than of $F_Q$.
\end{remark}

%% ============================================================
\section{Decoherence as Dimensional Projection}
\label{sec:decoherence}
We now apply the information-geometric framework of Section~\ref{sec:geometry} to the decoherence process. The key question is this: when a quantum system decoheres---when environmental interactions destroy coherences in a preferred pointer basis---what happens to the information-geometric structure of the state manifold? We show that decoherence acts as a geometric projection: it annihilates the $d^2 - d$ off-diagonal directions of the quantum state manifold, collapsing it onto the $(d-1)$-dimensional classical probability simplex. The quantum Fisher information metric restricts, under this projection, to the classical Fisher--Rao metric. This is a precise, quantitative account of the quantum-to-classical transition in information-geometric terms.

\begin{definition}[Decoherence channel]
\label{def:decoherence-channel}
Fix an orthonormal pointer basis $\{|k\rangle\}_{k=1}^d$. The \emph{complete dephasing channel} is the CPTP map
\begin{equation}
\label{eq:dephasing}
\Lambda_{\mathrm{dec}}(\rho) = \sum_{k=1}^d |k\rangle\langle k|\, \rho\, |k\rangle\langle k|.
\end{equation}
This sets all off-diagonal elements to zero: $[\Lambda_{\mathrm{dec}}(\rho)]_{kl} = \delta_{kl}\, \langle k|\rho|k\rangle$. The Kraus operators are the rank-one projectors $K_k = |k\rangle\langle k|$, satisfying $\sum_k K_k^\dagger K_k = I$.
\end{definition}

\begin{theorem}[Dimensional collapse]
\label{thm:dim-collapse}
Let $\mathcal{M}_Q$ be the quantum state manifold of a $d$-dimensional system $(\dim \mathcal{M}_Q = d^2 - 1)$. The decoherence channel $\Lambda_{\mathrm{dec}}$ maps $\mathcal{M}_Q$ onto the classical probability simplex $\Delta_{d-1}$ $(\dim \Delta_{d-1} = d-1)$. Specifically:

\begin{enumerate}
\item[\emph{(i)}] The image $\Lambda_{\mathrm{dec}}(\mathcal{M}_Q) = \Delta_{d-1}$.

\item[\emph{(ii)}] The $d^2 - d$ off-diagonal parameters become unestimable: for any smooth parameterization $\theta = (\theta_{\mathrm{diag}}, \theta_{\mathrm{off}})$ of $\mathcal{M}_Q$ in which $\theta_{\mathrm{diag}}$ parameterizes the diagonal elements and $\theta_{\mathrm{off}}$ the off-diagonal elements, the QFIM of $\Lambda_{\mathrm{dec}}[\rho(\theta)]$ satisfies
\begin{equation}
\label{eq:offdiag-vanish}
[F_Q(\Lambda_{\mathrm{dec}}[\rho])]_{\alpha\beta} = 0 \quad \text{whenever } \alpha \in \theta_{\mathrm{off}} \text{ or } \beta \in \theta_{\mathrm{off}}.
\end{equation}
In particular, $\mathrm{rank}\, F_Q(\Lambda_{\mathrm{dec}}[\rho]) \leq d - 1$.

\item[\emph{(iii)}] The dimensional collapse ratio is
\begin{equation}
\label{eq:collapse-ratio}
\frac{\dim \Delta_{d-1}}{\dim \mathcal{M}_Q} = \frac{d-1}{d^2-1} = \frac{1}{d+1},
\end{equation}
which vanishes as $d \to \infty$.
\end{enumerate}
\end{theorem}

\begin{proof}
Part~(i) is immediate: $\Lambda_{\mathrm{dec}}$ sends $\rho$ to $\mathrm{diag}(\langle k|\rho|k\rangle)$, which is a point in $\Delta_{d-1}$, and every point in $\Delta_{d-1}$ is the image of some diagonal state.  Part~(iii) follows directly from $(d-1)/(d^2-1) = 1/(d+1)$.

For Part~(ii): the decohered state $\rho_{\mathrm{dec}}$ depends only on $\theta_{\mathrm{diag}}$, so $\partial \rho_{\mathrm{dec}}/\partial \theta_{\mathrm{off}}^\alpha = 0$, which gives $L_\alpha = 0$ (or zero on the support of $\rho_{\mathrm{dec}}$) and hence $[F_Q]_{\alpha\beta} = 0$ whenever $\alpha$ or $\beta$ is off-diagonal.  The surviving diagonal block reduces to the classical Fisher information:
\begin{equation}
\label{eq:qfim-classical}
[F_Q]_{ij}^{\mathrm{diag}} = \sum_{k=1}^d \frac{1}{p_k}\, \frac{\partial p_k}{\partial \theta_{\mathrm{diag}}^i}\, \frac{\partial p_k}{\partial \theta_{\mathrm{diag}}^j} = [F_{\mathrm{cl}}]_{ij},
\end{equation}
which has rank at most $d - 1$.
\end{proof}

\begin{theorem}[Metric restriction]
\label{thm:metric-restriction}
Under $\Lambda_{\mathrm{dec}}$, the Bures metric on $\mathcal{M}_Q$ restricts to the Fisher--Rao metric on $\Delta_{d-1}$:
\begin{equation}
\label{eq:metric-restriction}
ds^2_{\mathrm{Bures}}\big|_{\Delta_{d-1}} = \frac{1}{4} \sum_{k=1}^d \frac{dp_k^2}{p_k} = \frac{1}{4}\, ds^2_{\mathrm{Fisher\text{-}Rao}}.
\end{equation}
The quantum state manifold with its Bures geometry is a $(d^2 - 1)$-dimensional extension of the classical simplex with its Fisher--Rao geometry. Decoherence is the projection that discards the quantum extension.
\end{theorem}

\begin{proof}
For diagonal states, the SLD is diagonal: $\langle k|L_i|k\rangle = (\partial_i p_k)/p_k$.  The Bures line element becomes $ds^2_{\mathrm{Bures}} = \frac{1}{4}\sum_k dp_k^2/p_k$, which is $\frac{1}{4}$ times the Fisher--Rao metric.  This identification is standard \cite{braunstein1994statistical}.
\end{proof}

\begin{theorem}[Orthogonal decomposition]
\label{thm:orthogonal-decomposition}
Let $\rho = \mathrm{diag}(p_1, \ldots, p_d)$ be a full-rank diagonal state.  The tangent space $T_\rho \mathcal{M}_Q$---the space of traceless Hermitian matrices at $\rho$---admits an orthogonal decomposition under the Bures metric:
\begin{equation}
\label{eq:orthogonal-decomposition}
T_\rho \mathcal{M}_Q = T_\rho^{\mathrm{diag}} \oplus T_\rho^{\mathrm{off}},
\end{equation}
where:
\begin{itemize}
\item $T_\rho^{\mathrm{diag}}$ is the $(d-1)$-dimensional traceless diagonal subspace,
\item $T_\rho^{\mathrm{off}}$ is the $(d^2-d)$-dimensional off-diagonal Hermitian subspace.
\end{itemize}
The differential $d\Lambda_{\mathrm{dec}}|_\rho = \Lambda_{\mathrm{dec}}$ acts as the identity on $T_\rho^{\mathrm{diag}}$ and as zero on $T_\rho^{\mathrm{off}}$.  Therefore $d\Lambda_{\mathrm{dec}}|_\rho$ is the $g_B$-orthogonal projector onto $T_\rho^{\mathrm{diag}}$.
\end{theorem}

\begin{proof}
We proceed in three steps: decomposition, orthogonality, and identification of the differential as a projector.

\emph{Step 1: Decomposition.}  At a diagonal state $\rho = \mathrm{diag}(p)$ with all $p_k > 0$, any tangent vector $X \in T_\rho \mathcal{M}_Q$ (a traceless Hermitian matrix) decomposes uniquely as $X = X_D + X_O$, where $X_D$ is the diagonal part and $X_O$ is the off-diagonal part.  Since $\mathrm{Tr}[X] = \mathrm{Tr}[X_D] = 0$, the diagonal part $X_D$ lies in the $(d-1)$-dimensional traceless diagonal subspace $T_\rho^{\mathrm{diag}}$, and the off-diagonal part $X_O$ lies in the $(d^2 - d)$-dimensional off-diagonal Hermitian subspace $T_\rho^{\mathrm{off}}$.

\emph{Step 2: Orthogonality under the Bures metric.}  The Bures inner product of two tangent vectors $A, B \in T_\rho \mathcal{M}_Q$ is
\[
g_B(A, B) = \frac{1}{4}\,\mathrm{Tr}\!\left[\rho\,\frac{L_A L_B + L_B L_A}{2}\right],
\]
where $L_A$ is the SLD satisfying $A = (\rho L_A + L_A \rho)/2$.

For a diagonal tangent vector $X_D$ with $(X_D)_{kk} = x_k$, the SLD is diagonal: $(L_{X_D})_{kk} = x_k / p_k$, with all off-diagonal entries zero.

For an off-diagonal tangent vector $X_O$ with $(X_O)_{mn} = z_{mn}$ ($m \neq n$), the SLD is off-diagonal: $(L_{X_O})_{mn} = 2\, z_{mn} / (p_m + p_n)$, with all diagonal entries zero.

The anticommutator $\{L_{X_D}, L_{X_O}\} = L_{X_D} L_{X_O} + L_{X_O} L_{X_D}$.  Since $L_{X_D}$ is diagonal and $L_{X_O}$ is off-diagonal, the product $L_{X_D} L_{X_O}$ is off-diagonal, and likewise $L_{X_O} L_{X_D}$.  Their sum $\{L_{X_D}, L_{X_O}\}$ is therefore off-diagonal.  Multiplying by the diagonal matrix $\rho$ yields a matrix that is still off-diagonal.  The trace of an off-diagonal matrix is zero:
\[
g_B(X_D, X_O) = \frac{1}{4}\,\mathrm{Tr}\!\left[\rho\,\frac{\{L_{X_D}, L_{X_O}\}}{2}\right] = 0.
\]
This establishes $T_\rho^{\mathrm{diag}} \perp_{g_B} T_\rho^{\mathrm{off}}$.

\emph{Step 3: The differential as orthogonal projector.}  The dephasing channel acts linearly on tangent vectors: $d\Lambda_{\mathrm{dec}}|_\rho(X) = \Lambda_{\mathrm{dec}}(X)$.  For a diagonal perturbation $X_D$, dephasing preserves it: $\Lambda_{\mathrm{dec}}(X_D) = X_D$.  For an off-diagonal perturbation $X_O$, dephasing annihilates it: $\Lambda_{\mathrm{dec}}(X_O) = 0$.  Combined with the Bures-orthogonality of $T_\rho^{\mathrm{diag}}$ and $T_\rho^{\mathrm{off}}$, this identifies $d\Lambda_{\mathrm{dec}}|_\rho$ as the $g_B$-orthogonal projector onto $T_\rho^{\mathrm{diag}}$.
\end{proof}

\begin{remark}[Geometric content]
\label{rem:geometric-content}
At full-rank diagonal states, the orthogonal decomposition reveals that classical and quantum degrees of freedom are geometrically perpendicular in the Bures metric.  The differential of the dephasing map does not merely reduce dimension---it is the Bures-orthogonal projector onto the classical tangent subspace.  This is the geometric content of the statement that decoherence preserves classical information while destroying quantum information: at these states, the two types of information live in orthogonal subspaces.  Theorem~\ref{thm:orthogonal-decomposition} reveals the local differential-geometric structure underlying the metric restriction of Theorem~\ref{thm:metric-restriction}.
\end{remark}

\begin{corollary}[$D_{\mathrm{eff}}$ collapse]
\label{cor:deff-collapse}
Under parameterization of the full state manifold $(p = d^2 - 1)$, the effective quantum dimensionality satisfies
\begin{equation}
\label{eq:deff-bound}
D_{\mathrm{eff}}^Q(\Lambda_{\mathrm{dec}}[\rho]) \leq d - 1,
\end{equation}
with equality when all $d - 1$ diagonal parameters are equally estimable (e.g., when $p_k = 1/d$ for all $k$).

In contrast, for a generic pure state with full coherence, $D_{\mathrm{eff}}^Q(\rho)$ can approach $d^2 - 1$.
\end{corollary}

\begin{proof}
By Theorem~\ref{thm:dim-collapse}(ii), the QFIM of $\Lambda_{\mathrm{dec}}[\rho]$ has at most $d - 1$ nonzero eigenvalues. Since $D_{\mathrm{eff}}^Q \leq \mathrm{rank}(F_Q)$ by the participation ratio bound (Definition~\ref{def:deff}), the inequality~\eqref{eq:deff-bound} follows.

For equality at the uniform distribution $p_k = 1/d$: by the permutation symmetry of the uniform distribution under relabeling of pointer states, the classical Fisher information matrix~\eqref{eq:qfim-classical} is proportional to the identity on the $(d-1)$-dimensional subspace orthogonal to the trace constraint. All $d - 1$ eigenvalues are therefore equal, and $D_{\mathrm{eff}}^Q = (d-1)^2 \cdot 1^2 / ((d-1) \cdot 1^2) = d - 1$.
\end{proof}

The $d^2 - d$ dimensions destroyed by decoherence are precisely the quantum coherences---the off-diagonal elements of $\rho$ that encode phase relationships between pointer states. These dimensions have no classical counterpart. Decoherence does not destroy information randomly; it eliminates specifically those information-geometric dimensions that distinguish quantum from classical state spaces.

The surviving $d - 1$ dimensions are the populations---the probabilities over pointer states that can be carried into the observable, classical world. This is the information-geometric content of Zurek's observation that ``pointer states survive decoherence'' \cite{zurek2003decoherence}: in our framework, they survive because they span the submanifold onto which the decoherence channel projects.

The collapse ratio $1/(d+1)$ reveals the severity of dimensional loss: for $d = 2$, one-third of the information-geometric dimensions survive; for $d = 1024$ (10 qubits), less than $0.1\%$. As $d$ grows, the classical world becomes a vanishingly thin projection of quantum state space.

\subsection{The classical--quantum split away from the simplex}
\label{sec:bures-angle}

Theorem~\ref{thm:orthogonal-decomposition} establishes orthogonality at diagonal states.  The main result of this paper answers the natural next question: what happens at states that carry coherence?  For qubits, a single principal angle characterizes the full relationship between the two subspaces, and we compute it exactly.

\begin{theorem}[Bures separation angle]
\label{thm:bures-angle}
Let $\rho = (I + \mathbf{r} \cdot \boldsymbol{\sigma})/2$ be a full-rank qubit state $(|\mathbf{r}| < 1)$ with Bloch vector $\mathbf{r} = (r_x, r_y, r_z)$ in the pointer basis $\{|0\rangle, |1\rangle\}$.  Write $r_\perp = \sqrt{r_x^2 + r_y^2}$ for the coherence magnitude and $r_z$ for the population imbalance.  The tangent space $T_\rho \mathcal{M}_Q$ decomposes as $T_\rho^{\mathrm{diag}} \oplus T_\rho^{\mathrm{off}}$ (as vector spaces), where $T_\rho^{\mathrm{diag}} = \mathrm{span}\{\sigma_z\}$ and $T_\rho^{\mathrm{off}} = \mathrm{span}\{\sigma_x, \sigma_y\}$.

The principal angle $\theta_B$ between $T_\rho^{\mathrm{diag}}$ and $T_\rho^{\mathrm{off}}$ under the Bures metric satisfies
\begin{equation}
\label{eq:bures-angle}
\cos^2 \theta_B = \frac{r_z^2\, r_\perp^2}{(1 - r_\perp^2)(1 - r_z^2)}.
\end{equation}
Equivalently,
\begin{equation}
\label{eq:bures-angle-sin}
\sin^2 \theta_B = \frac{1 - |\mathbf{r}|^2}{(1 - r_\perp^2)(1 - r_z^2)}.
\end{equation}
\end{theorem}

\begin{proof}
The Bures metric on qubit state space in Bloch coordinates is
\begin{equation}
\label{eq:bures-bloch}
g_{ij}^B = \frac{1}{4}\!\left[\delta_{ij} + \frac{r_i\, r_j}{1 - |\mathbf{r}|^2}\right].
\end{equation}
This follows from the QFIM $[F_Q]_{ij} = \delta_{ij} + r_i r_j/(1-|\mathbf{r}|^2)$, which is obtained by computing the diagonal contribution $\hat{r}_i \hat{r}_j/(1-|\mathbf{r}|^2)$ and the off-diagonal contribution $\delta_{ij} - \hat{r}_i \hat{r}_j$ from the SLD equation~\eqref{eq:sld-matrix} in the eigenbasis of~$\rho$ (see Appendix, or \cite{braunstein1994statistical}).

The principal angle between $T_\rho^{\mathrm{diag}}$ and $T_\rho^{\mathrm{off}}$ is found by maximizing $g_B(\hat{e}_z, v)/(\lVert\hat{e}_z\rVert \cdot \lVert v\rVert)$ over $v \in T_\rho^{\mathrm{off}}$.  The maximizer is $v_\perp = (r_x, r_y, 0)/r_\perp$ (the projection of~$\mathbf{r}$ onto the equatorial plane), since this is the unique direction in $T_\rho^{\mathrm{off}}$ with nonvanishing inner product against~$\hat{e}_z$.  Using~\eqref{eq:bures-bloch}:
\begin{align*}
g_B(\hat{e}_z, v_\perp) &= \frac{r_z\, r_\perp}{4(1 - |\mathbf{r}|^2)}, \\[3pt]
g_B(\hat{e}_z, \hat{e}_z) &= \frac{1 - r_\perp^2}{4(1 - |\mathbf{r}|^2)}, \\[3pt]
g_B(v_\perp, v_\perp) &= \frac{1 - r_z^2}{4(1 - |\mathbf{r}|^2)}.
\end{align*}
The squared cosine of the principal angle is
\[
\cos^2\theta_B = \frac{[g_B(\hat{e}_z, v_\perp)]^2}{g_B(\hat{e}_z, \hat{e}_z) \cdot g_B(v_\perp, v_\perp)} = \frac{r_z^2\, r_\perp^2}{(1 - r_\perp^2)(1 - r_z^2)}.
\]
The sine formula follows from $\sin^2\theta_B = 1 - \cos^2\theta_B$ and the identity $(1 - r_\perp^2)(1 - r_z^2) - r_z^2 r_\perp^2 = 1 - r_\perp^2 - r_z^2 = 1 - |\mathbf{r}|^2$.
\end{proof}

\begin{corollary}[Structure of the classical--quantum split]
\label{cor:angle-structure}
\quad
\begin{enumerate}
\item[\emph{(i)}] $\theta_B = \pi/2$ if and only if $r_\perp = 0$ \emph{or} $r_z = 0$.  In particular, the subspaces are Bures-orthogonal at all diagonal states (recovering Theorem~\ref{thm:orthogonal-decomposition}) and at all equatorial states (where the population imbalance vanishes).
\item[\emph{(ii)}] $\theta_B < \pi/2$ if and only if both $r_\perp \neq 0$ and $r_z \neq 0$, i.e., the state simultaneously carries nonzero coherence and nonzero population imbalance.
\item[\emph{(iii)}] For fixed $r_z \neq 0$, the dephasing map $r_\perp \mapsto (1-2\gamma)\, r_\perp$ monotonically increases $\theta_B$ toward $\pi/2$ as $\gamma \to 1/2$.
\end{enumerate}
\end{corollary}

\begin{proof}
Parts (i) and (ii) follow directly from the formula~\eqref{eq:bures-angle}: $\cos^2\theta_B = 0$ if and only if $r_z r_\perp = 0$.  For part~(iii), under dephasing with parameter~$\gamma$, $r_\perp \mapsto (1-2\gamma) r_\perp$ while $r_z$ is preserved.  Substituting into~\eqref{eq:bures-angle}:
\[
\cos^2\theta_B(\gamma) = \frac{r_z^2 (1-2\gamma)^2 r_\perp^2}{(1 - (1-2\gamma)^2 r_\perp^2)(1 - r_z^2)}.
\]
The numerator decreases monotonically in $\gamma$ while the denominator increases (or stays constant), so $\cos^2\theta_B$ decreases monotonically and $\theta_B$ increases toward~$\pi/2$.
\end{proof}

\begin{remark}[Geometric interpretation]
\label{rem:angle-interpretation}
The classical--quantum tangent split is oblique exactly when the state simultaneously carries nonzero coherence and nonzero population imbalance.  The ``mixing term'' $\cos^2\theta_B$ is controlled by the product $|r_z|\, r_\perp$: population imbalance alone (diagonal states) or coherence alone (equatorial states) each preserve orthogonality, but their combination breaks it.

As dephasing eliminates coherence ($r_\perp \to 0$) while preserving populations ($r_z$ fixed), the angle increases monotonically to~$\pi/2$, realizing the orthogonal decomposition of Theorem~\ref{thm:orthogonal-decomposition}.  The metric separation between classical and quantum degrees of freedom is state-dependent and is sharpened by dephasing: at full dephasing, the split is perfectly orthogonal.  Figure~\ref{fig:bures-angle} illustrates this structure across the interior of the Bloch disk.
\end{remark}

\begin{corollary}[Estimation-theoretic interpretation]
\label{cor:estimation}
For a full-rank qubit state $\rho$ with Bloch vector~$\mathbf{r}$ $(|\mathbf{r}| < 1)$, the quantum Fisher information for the population parameter~$r_z$ in the Bloch parameterization is $[F_Q]_{zz} = (1 - r_\perp^2)/(1 - |\mathbf{r}|^2)$.  When the coherence parameters $(r_x, r_y)$ are treated as nuisance parameters and eliminated via the Schur complement \cite{zhang2025hybrid,albarelli2019restoring}, the effective Fisher information for~$r_z$ is
\begin{equation}
\label{eq:fzz-eff}
[F_Q]_{zz}^{\mathrm{eff}} \;=\; [F_Q]_{zz}\,\sin^2\theta_B \;=\; \frac{1}{1 - r_z^2}.
\end{equation}
The nuisance penalty $[F_Q]_{zz}\cos^2\theta_B$ quantifies the fraction of population Fisher information that is entangled with coherence parameters.  The Bures separation angle determines the fraction of the raw SLD Fisher information for~$r_z$ that survives elimination of the coherence nuisance parameters.  (This is a bound on estimation precision via the multiparameter quantum Cram\'er--Rao inequality; whether the bound is attainable depends on the compatibility of optimal measurements for the different parameters \cite{albarelli2019restoring,wang2024multiparameter}.)
\end{corollary}

\begin{proof}
The full QFIM in Bloch coordinates is $[F_Q]_{ij} = \delta_{ij} + r_i r_j/(1 - |\mathbf{r}|^2)$, which is a rank-one update of the identity: $F_Q = I + \mathbf{r}\mathbf{r}^T/(1-|\mathbf{r}|^2)$.  By the Sherman--Morrison formula,
\[
(F_Q^{-1})_{ij} \;=\; \delta_{ij} - r_i r_j,
\]
so $(F_Q^{-1})_{zz} = 1 - r_z^2$.  For the partitioned QFIM with $r_z$ as the parameter of interest and $(r_x, r_y)$ as nuisance parameters, the Schur complement gives $[F_Q]_{zz}^{\mathrm{eff}} = 1/(F_Q^{-1})_{zz} = 1/(1 - r_z^2)$.  Comparing with $[F_Q]_{zz} = (1 - r_\perp^2)/(1 - |\mathbf{r}|^2)$:
\[
\frac{[F_Q]_{zz}^{\mathrm{eff}}}{[F_Q]_{zz}} = \frac{1 - |\mathbf{r}|^2}{(1 - r_\perp^2)(1 - r_z^2)} = \sin^2\theta_B
\]
by Theorem~\ref{thm:bures-angle}.
\end{proof}

\begin{figure}[H]
\centering
\includegraphics[width=0.7\textwidth]{fig5_bures_angle.pdf}
\caption{\textbf{Bures separation angle across the Bloch disk.}  The principal angle $\theta_B$ between the diagonal subspace $T_\rho^{\mathrm{diag}}$ and the off-diagonal subspace $T_\rho^{\mathrm{off}}$, plotted on the $(r_\perp, r_z)$ cross-section of the Bloch ball interior $(|\mathbf{r}| < 1)$.  The angle equals $\pi/2$ (perfect orthogonality) on both axes ($r_\perp = 0$ and $r_z = 0$) and decreases toward the pure-state boundary where both coherence and population imbalance are large.  Dashed horizontal arrows show dephasing trajectories at fixed $r_z$, along which $\theta_B$ increases monotonically toward~$\pi/2$.}
\label{fig:bures-angle}
\end{figure}

\subsection{Generalization to qudits}
\label{sec:qudit-angles}

Theorem~\ref{thm:bures-angle} gives the exact principal angle for qubits.  We now characterize the classical--quantum tangent split for general $d$-dimensional systems.  Since $\dim T_\rho^{\mathrm{diag}} = d - 1 \leq d^2 - d = \dim T_\rho^{\mathrm{off}}$, there are $d - 1$ principal angles $\theta_1, \ldots, \theta_{d-1}$ between the two subspaces (defined as the singular-value angles of the normalized cross-Gram matrix).

\begin{theorem}[Qudit principal angles]
\label{thm:qudit-angles}
Let $\rho$ be a full-rank $d$-dimensional state.  Write $\rho = \mathrm{diag}(p) + C$, where $\mathrm{diag}(p)$ has entries $p_k = \langle k|\rho|k\rangle > 0$ and $C$ collects the off-diagonal entries $C_{kl} = \langle k|\rho|l\rangle$ ($k \neq l$).
\begin{enumerate}
\item[\emph{(i)}] If $\rho$ is diagonal ($C = 0$), all $d - 1$ principal angles between $T_\rho^{\mathrm{diag}}$ and $T_\rho^{\mathrm{off}}$ equal $\pi/2$ \emph{(Theorem~\ref{thm:orthogonal-decomposition})}.
\item[\emph{(ii)}] \emph{(Near-simplex expansion.)} For states near the classical simplex ($\lVert C \rVert$ small), the cross-Gram matrix has entries
\begin{equation}
\label{eq:cross-gram}
g_B(D_a, R_{kl}) = -\frac{\mathrm{Re}[C_{kl}]}{2(p_k + p_l)}\!\left(\frac{x_a^k}{p_k} + \frac{x_a^l}{p_l}\right) + O(\lVert C \rVert^2),
\end{equation}
where $D_a = \mathrm{diag}(x_a^1, \ldots, x_a^d)$ is a traceless diagonal tangent vector and $R_{kl} = |k\rangle\langle l| + |l\rangle\langle k|$ is a real off-diagonal tangent vector.
\item[\emph{(iii)}] For $d \geq 3$: if $\lVert C \rVert$ is sufficiently small and $C \neq 0$, at least one principal angle is strictly less than~$\pi/2$.  (For $d = 2$, the exact criterion is given by Theorem~\ref{thm:bures-angle}: $\theta_B < \pi/2$ iff both $r_\perp \neq 0$ and $r_z \neq 0$.)
\item[\emph{(iv)}] Under partial dephasing $C_{kl} \mapsto (1-2\gamma)\, C_{kl}$, the cross-Gram entries scale as $(1-2\gamma)$ at leading order, and all principal angles converge to $\pi/2$ as $\gamma \to 1/2$.
\end{enumerate}
\end{theorem}

\begin{proof}
\emph{Part (i), ``$\Leftarrow$'':} This is Theorem~\ref{thm:orthogonal-decomposition}.

\emph{Part (ii):}  Write $\rho = \mathrm{diag}(p) + C$ with $\lVert C \rVert$ small.  The SLD $L_a$ for the diagonal tangent vector $D_a$ satisfies $\rho L_a + L_a \rho = 2 D_a$.  Expanding $L_a = L_a^{(0)} + L_a^{(1)} + O(\lVert C \rVert^2)$:

At zeroth order, $(L_a^{(0)})_{kk} = x_a^k / p_k$ with all off-diagonal entries zero.  At first order, the Lyapunov equation gives $(L_a^{(1)})_{mn} = -C_{mn}(x_a^m/p_m + x_a^n/p_n)/(p_m + p_n)$ for $m \neq n$ and $(L_a^{(1)})_{kk} = 0$.

For $R_{kl}$, the zeroth-order SLD is $(L_R^{(0)})_{kl} = (L_R^{(0)})_{lk} = 2/(p_k + p_l)$, all other entries zero.  Its first-order diagonal corrections are $(L_R^{(1)})_{kk} = -2\,\mathrm{Re}[C_{kl}]/[p_k(p_k + p_l)]$ and $(L_R^{(1)})_{ll} = -2\,\mathrm{Re}[C_{kl}]/[p_l(p_k + p_l)]$.

The Bures inner product $g_B(D_a, R_{kl}) = \frac{1}{4}\,\mathrm{Re}\,\mathrm{Tr}[\rho\, L_a L_R]$ receives three first-order contributions:
\begin{itemize}
\item $C$ acting on zeroth-order SLDs: $+\frac{\mathrm{Re}[C_{kl}]}{2(p_k+p_l)}\!\left(\frac{x_a^k}{p_k} + \frac{x_a^l}{p_l}\right)$.
\item $L_a^{(1)}$ with $L_R^{(0)}$: $-\frac{\mathrm{Re}[C_{kl}]}{2(p_k+p_l)}\!\left(\frac{x_a^k}{p_k} + \frac{x_a^l}{p_l}\right)$.
\item $L_R^{(1)}$ with $L_a^{(0)}$: $-\frac{\mathrm{Re}[C_{kl}]}{2(p_k+p_l)}\!\left(\frac{x_a^k}{p_k} + \frac{x_a^l}{p_l}\right)$.
\end{itemize}
Summing yields~\eqref{eq:cross-gram}.  (The analogous formula for the imaginary direction $I_{kl} = i|k\rangle\langle l| - i|l\rangle\langle k|$ replaces $\mathrm{Re}[C_{kl}]$ with $\mathrm{Im}[C_{kl}]$.)

\emph{Part (iii):}  Suppose $C_{kl} \neq 0$ for some $k < l$ with $d \geq 3$.  Choose the traceless diagonal vector $D_a$ with $x_a^k = 1$, $x_a^m = -1$ for a single $m \notin \{k, l\}$, and $x_a^j = 0$ otherwise.  Then $x_a^k/p_k + x_a^l/p_l = 1/p_k \neq 0$, and~\eqref{eq:cross-gram} gives $g_B(D_a, R_{kl}) = -\mathrm{Re}[C_{kl}]/[2\,p_k(p_k + p_l)] + O(\lVert C \rVert^2)$.  For sufficiently small $\lVert C \rVert$, this is nonzero (using $I_{kl}$ instead if $\mathrm{Re}[C_{kl}] = 0$ but $\mathrm{Im}[C_{kl}] \neq 0$).  The key point is that for $d \geq 3$, the $(d-1)$-dimensional space of traceless diagonal vectors always contains a direction with $x_a^k/p_k + x_a^l/p_l \neq 0$, because $x_a^k$ and $x_a^l$ can be varied independently (via third levels $m \notin \{k,l\}$).  For $d = 2$, the unique traceless direction $D_1 \propto \mathrm{diag}(1,-1)$ satisfies $x_1^1/p_1 + x_1^2/p_2 = 1/p_1 - 1/p_2 = 0$ at equal populations, allowing orthogonality to persist; the exact qubit criterion is given by Theorem~\ref{thm:bures-angle}.

\emph{Part (iv):}  Under $C \mapsto (1-2\gamma)\, C$, equation~\eqref{eq:cross-gram} scales the cross-Gram as $(1-2\gamma)$ at leading order.  As $\gamma \to 1/2$, the state converges to $\mathrm{diag}(p)$ and all cross-Gram entries vanish, so all principal angles converge to~$\pi/2$ by continuity.
\end{proof}

\begin{remark}[Qubit vs.\ qudit orthogonality]
\label{rem:qubit-qudit}
The orthogonality criterion sharpens from $d = 2$ to $d \geq 3$.  For qubits, equal populations ($r_z = 0$) suffice for orthogonality even in the presence of coherence---the equatorial plane of the Bloch ball has $\theta_B = \pi/2$.  For $d \geq 3$, Part~(iii) shows that near the classical simplex, any nonzero coherence breaks orthogonality.  The structural reason is dimensional: for $d = 2$, the traceless diagonal subspace is one-dimensional, and the single direction $\mathrm{diag}(1,-1)$ has $x^1/p_1 + x^2/p_2 = 0$ at equal populations.  For $d \geq 3$, the $(d-1)$-dimensional traceless diagonal space always contains directions that couple to any given off-diagonal entry.  Theorem~\ref{thm:spectator-orthogonality} below shows that this is not merely a local phenomenon: for $d \geq 3$, any nonzero coherence with a decoupled spectator breaks orthogonality.  We conjecture (Conjecture~\ref{conj:global-orthogonality}) that only diagonal states achieve orthogonality; a numerical survey (Figure~\ref{fig:qudit-angles}) supports this for $d = 3, 4, 5$.
\end{remark}

\begin{figure}[H]
\centering
\includegraphics[width=\textwidth]{fig6_qudit_angles.pdf}
\caption{\textbf{Numerical survey of qudit principal angles.}  For each dimension $d \in \{3, 4, 5\}$, $2{,}000$ random full-rank density matrices are drawn from the Hilbert--Schmidt measure (blue histogram), and the minimum principal angle $\min_a \theta_a$ between the diagonal and off-diagonal tangent subspaces is computed via the Bures metric.  The dashed red line marks $90^\circ$ (perfect orthogonality).  The solid orange line marks the best angle found by adversarial optimization ($100$ L-BFGS-B starts per dimension; see Appendix~\ref{app:numerics}).  No state---random or optimized---approaches orthogonality; the distribution shifts further from $90^\circ$ as $d$ increases, supporting Conjecture~\ref{conj:global-orthogonality}.}
\label{fig:qudit-angles}
\end{figure}

\begin{theorem}[Spectator orthogonality criterion for $d \geq 3$]
\label{thm:spectator-orthogonality}
Let $d \geq 3$ and let $\rho$ be a full-rank state with $\rho_{st} \neq 0$ for some $s \neq t$.  If there exists a spectator index $r \notin \{s,t\}$ such that $\rho_{rs} = \rho_{rt} = 0$, then the Bures-orthogonal decomposition $T_\rho^{\mathrm{diag}} \perp_{g_B} T_\rho^{\mathrm{off}}$ fails.  In particular, orthogonality fails whenever the pointer-basis coherence graph---with vertices $\{1,\ldots,d\}$ and edges $\{(k,l) : \rho_{kl} \neq 0\}$---contains an edge whose endpoints share a common non-neighbor.
\end{theorem}

\begin{proof}
We prove that if $\rho_{st} \neq 0$ and some spectator $r$ has $\rho_{rs} = \rho_{rt} = 0$, then the cross-Gram matrix $G_{DO}$ has a nonzero entry.

\emph{Step 1: Reduction to SLD diagonal entries.}
For any diagonal $X_D = \mathrm{diag}(x_1, \ldots, x_d)$ and off-diagonal Hermitian~$X_O$, the standard SLD metric identity gives
\begin{equation}
\label{eq:sld-identity}
g_B(X_D, X_O) = \tfrac{1}{2}\,\mathrm{Re}\,\mathrm{Tr}[X_D\, L_{X_O}] = \tfrac{1}{2}\sum_k x_k\, (L_{X_O})_{kk},
\end{equation}
where $L_{X_O}$ is the unique Hermitian SLD satisfying $\rho L_{X_O} + L_{X_O} \rho = 2 X_O$, and $(L_{X_O})_{kk} \in \mathbb{R}$ by Hermiticity.  (The identity follows from $\mathrm{Tr}[\rho\,\{L_X, L_Y\}] = 2\,\mathrm{Tr}[X\, L_Y]$, which is obtained by substituting $\rho L_X = 2X - L_X\rho$ into the left side and applying cyclicity of the trace.)  Since $X_D$ is diagonal, only the pointer-basis diagonal entries of~$L_{X_O}$ contribute. Equation~\eqref{eq:sld-identity} vanishes for all traceless $(x_1, \ldots, x_d)$ if and only if $(L_{X_O})_{kk}$ is independent of~$k$.  Thus, Bures-orthogonality requires every off-diagonal SLD to have a flat diagonal in the pointer basis.

\emph{Step 2: Diagonal extraction from the Lyapunov equation.}
The $(k,k)$ entry of $\rho L + L\rho = 2X_O$ gives (using $(X_O)_{kk} = 0$):
\begin{equation}
\label{eq:lyapunov-diag}
p_k\,(L_{X_O})_{kk} = -\mathrm{Re}\!\left[\sum_{j \neq k} \rho_{kj}\,(L_{X_O})_{jk}\right].
\end{equation}

\emph{Step 3: Spectator decoupling.}
Consider the two off-diagonal tangent vectors $R_{st} = |s\rangle\langle t| + |t\rangle\langle s|$ and $I_{st} = i\,|s\rangle\langle t| - i\,|t\rangle\langle s|$.  Since $\rho_{rs} = \rho_{rt} = 0$, the $(r,s)$ and $(r,t)$ entries of the Lyapunov equation for any perturbation $X_O$ supported on the $(s,t)$ block form a homogeneous $2 \times 2$ system for $((L_{X_O})_{rs}, (L_{X_O})_{rt})$ with determinant $(p_r + p_s)(p_r + p_t) - |\rho_{st}|^2 > 0$.  Therefore $(L_{X_O})_{rs} = (L_{X_O})_{rt} = 0$, and~\eqref{eq:lyapunov-diag} with $k = r$ gives
\[
(L_{X_O})_{rr} = 0 \quad \text{for both } X_O = R_{st} \text{ and } X_O = I_{st}.
\]
For the diagonal entries at~$s$, equation~\eqref{eq:lyapunov-diag} with $k = s$ (and $\rho_{sr} = 0$) gives
\begin{equation}
\label{eq:Lss-formula}
(L_{X_O})_{ss} = -\mathrm{Re}[\rho_{st}\,(L_{X_O})_{ts}] / p_s.
\end{equation}
The $(s,t)$ entry of the Lyapunov equation becomes (with $\rho_{sr} = \rho_{tr} = 0$):
\begin{equation}
\label{eq:st-block}
(p_s + p_t)\,(L_{X_O})_{st} + \rho_{st}\bigl[(L_{X_O})_{tt} + (L_{X_O})_{ss}\bigr] = 2\,(X_O)_{st}.
\end{equation}
Both diagonal entries satisfy~\eqref{eq:Lss-formula} (with $s \leftrightarrow t$) with the same numerator $\mathrm{Re}[\rho_{st}\,(L_{X_O})_{ts}]$, so $(L_{X_O})_{tt} + (L_{X_O})_{ss} = -\mathrm{Re}[\rho_{st}\,(L_{X_O})_{ts}]\,(p_s^{-1} + p_t^{-1})$.  Substituting into~\eqref{eq:st-block} yields a scalar equation for $(L_{X_O})_{st}$ whose unique solution satisfies $(L_{X_O})_{st} \neq 0$ (since the right-hand side of~\eqref{eq:st-block} is nonzero).

Write $\rho_{st} = |\rho_{st}|\, e^{i\alpha}$.  Solving the resulting scalar equation for $(L_{X_O})_{st}$ yields the closed-form identities:
\begin{align}
\label{eq:Rformula}
\mathrm{Re}[\rho_{st}\,(L_{R_{st}})_{ts}] &= \frac{2\,|\rho_{st}|\,\cos\alpha}{(p_s + p_t)\bigl(1 - |\rho_{st}|^2/(p_s\,p_t)\bigr)}, \\[3pt]
\label{eq:Iformula}
\mathrm{Re}[\rho_{st}\,(L_{I_{st}})_{ts}] &= \frac{2\,|\rho_{st}|\,\sin\alpha}{(p_s + p_t)\bigl(1 - |\rho_{st}|^2/(p_s\,p_t)\bigr)},
\end{align}
where the common denominator is positive by the full-rank condition $|\rho_{st}|^2 < p_s\, p_t$.  Since $|\rho_{st}| > 0$, at least one of $\cos\alpha$, $\sin\alpha$ is nonzero; hence at least one of $(L_{R_{st}})_{ss}$ or $(L_{I_{st}})_{ss}$ is nonzero by~\eqref{eq:Lss-formula}.  Since the corresponding $(L_{X_O})_{rr} = 0$, the SLD diagonal is not flat: orthogonality fails.
\end{proof}

\begin{conjecture}[Global orthogonality criterion for $d \geq 3$]
\label{conj:global-orthogonality}
For $d \geq 3$, a full-rank state $\rho$ achieves the Bures-orthogonal decomposition $T_\rho^{\mathrm{diag}} \perp_{g_B} T_\rho^{\mathrm{off}}$ if and only if $\rho$ is diagonal in the pointer basis.
\end{conjecture}

Theorem~\ref{thm:spectator-orthogonality} proves this conjecture whenever the coherence graph contains an edge whose endpoints share a common non-neighbor.  The conjecture extends to fully-coherent states (complete coherence graph), where no spectator decouples.

Two rigorous results support the extension: Theorem~\ref{thm:qudit-angles}(iii) proves the conjecture in a neighborhood of the classical simplex for all~$d$, and a numerical survey of $2{,}000$ random full-rank states per dimension ($d = 3, 4, 5$, Figure~\ref{fig:qudit-angles}) finds no non-diagonal state achieving orthogonality.

\begin{remark}[Heuristic support for Conjecture~\ref{conj:global-orthogonality}]
Two additional observations, while not rigorous, suggest the conjecture holds globally.  First, the spectator formulas~\eqref{eq:Rformula}--\eqref{eq:Iformula} show that the failure of orthogonality is proportional to $|\rho_{st}|$; a discontinuous restoration of orthogonality when additional off-diagonal entries become nonzero would require exact cancellation among the SLD diagonal contributions.  Second, for $d = 3$ with all off-diagonal entries nonzero, the flat-diagonal condition yields an overdetermined algebraic system (twelve real constraints on three real degrees of freedom of the diagonalizing unitary); we have verified numerically for $50{,}000$ random fully-coherent qutrit states that no non-diagonal solution exists.  A complete proof would require showing these constraints are independent for all~$d$.
\end{remark}

%% ============================================================
\section{Physical Implications}
\label{sec:implications}

\subsection{Thermodynamic cost of dimensional collapse}
\label{sec:thermodynamics}
The dimensional collapse changes the non-equilibrium free energy of the state \cite{lostaglio2015description,korzekwa2016extraction}.

\begin{proposition}[Free energy of coherence]
\label{prop:free-energy-coherence}
Let $\rho$ be a quantum state with system Hamiltonian $H$, and let $\Lambda_{\mathrm{dec}}$ be the complete dephasing channel \emph{(Definition~\ref{def:decoherence-channel})}.  The free energy change under dephasing is
\begin{equation}
\label{eq:free-energy-general}
F(\rho) - F(\Lambda_{\mathrm{dec}}[\rho]) = \mathrm{Tr}\bigl[(\rho - \Lambda_{\mathrm{dec}}[\rho])\, H\bigr] + k_B T\, C_r(\rho),
\end{equation}
where $F(\rho) = \mathrm{Tr}[\rho H] - k_B T\, S(\rho)$ is the non-equilibrium free energy, $S(\rho) = -\mathrm{Tr}[\rho \ln \rho]$ is the von Neumann entropy, and $C_r(\rho) = S(\Lambda_{\mathrm{dec}}[\rho]) - S(\rho) \geq 0$ is the relative entropy of coherence \cite{streltsov2017colloquium}.
\end{proposition}

\begin{proof}
By definition of the non-equilibrium free energy:
\[
F(\rho) - F(\Lambda_{\mathrm{dec}}[\rho]) = \mathrm{Tr}[\rho H] - k_B T\, S(\rho) - \mathrm{Tr}[\Lambda_{\mathrm{dec}}[\rho]\, H] + k_B T\, S(\Lambda_{\mathrm{dec}}[\rho]).
\]
Rearranging:
\[
= \mathrm{Tr}\bigl[(\rho - \Lambda_{\mathrm{dec}}[\rho])\, H\bigr] + k_B T \bigl[S(\Lambda_{\mathrm{dec}}[\rho]) - S(\rho)\bigr].
\]
The entropy difference $S(\Lambda_{\mathrm{dec}}[\rho]) - S(\rho) \geq 0$ because $\Lambda_{\mathrm{dec}}$ is unital, so the eigenvalues of $\rho$ majorize those of $\Lambda_{\mathrm{dec}}[\rho]$ and von Neumann entropy is Schur-concave.  The identification $C_r(\rho) = S(\Lambda_{\mathrm{dec}}[\rho]) - S(\rho) = S(\rho \,\|\, \Lambda_{\mathrm{dec}}[\rho])$ is standard \cite{streltsov2017colloquium}: since $\Lambda_{\mathrm{dec}}[\rho] = \sum_k p_k |k\rangle\langle k|$ is diagonal in the pointer basis, $\mathrm{Tr}[\rho \ln \Lambda_{\mathrm{dec}}[\rho]] = \sum_k p_k \ln p_k = -S(\Lambda_{\mathrm{dec}}[\rho])$.
\end{proof}

\begin{corollary}[Dephasing-diagonal Hamiltonian]
\label{cor:diagonal-hamiltonian}
When $H$ is diagonal in the dephasing basis (i.e., every dephasing projector $|k\rangle\langle k|$ commutes with $H$), the energy term vanishes and
\begin{equation}
\label{eq:free-energy-diagonal-H}
F(\rho) - F(\Lambda_{\mathrm{dec}}[\rho]) = k_B T\, C_r(\rho) \;\geq\; 0.
\end{equation}
\end{corollary}

\begin{proof}
When $H$ is diagonal in the pointer basis, $H = \sum_k E_k |k\rangle\langle k|$.  Then $\mathrm{Tr}[\rho H] = \sum_k E_k \langle k|\rho|k\rangle = \sum_k E_k\, p_k = \mathrm{Tr}[\Lambda_{\mathrm{dec}}[\rho]\, H]$, so the energy term in~\eqref{eq:free-energy-general} vanishes.  Non-negativity follows from $C_r(\rho) \geq 0$ by unitality of $\Lambda_{\mathrm{dec}}$.
\end{proof}

\begin{remark}[Geometric correspondence]
\label{rem:geometric-correspondence}
The classical free energy $F(\Lambda_{\mathrm{dec}}[\rho])$ is associated with the $d - 1$ surviving population dimensions (Theorem~\ref{thm:dim-collapse}).  The coherent free energy $k_B T \cdot C_r(\rho)$ is associated with the $d^2 - d$ collapsing off-diagonal directions.  When $H$ is diagonal in the dephasing basis, the free energy decomposes cleanly along the orthogonal subspaces of Theorem~\ref{thm:orthogonal-decomposition}: the classical component on $T_\rho^{\mathrm{diag}}$ and the coherent component on $T_\rho^{\mathrm{off}}$.
\end{remark}

Proposition~\ref{prop:free-energy-coherence} decomposes the free energy change under dephasing into an energy shift and an entropic term $k_B T \cdot C_r$; when $H$ is diagonal in the dephasing basis, the energy shift vanishes and the free energy change equals $k_B T \cdot C_r$ exactly (Corollary~\ref{cor:diagonal-hamiltonian}).  A state with no coherence ($C_r = 0$) is already classical---no free energy is stored in coherences.  A maximally coherent state (e.g., $|\psi\rangle = d^{-1/2} \sum_k |k\rangle$) stores the maximum coherent free energy: $C_r = \ln d$.  The dimensions annihilated in Theorem~\ref{thm:dim-collapse}---the $d^2 - d$ off-diagonal directions whose Fisher information vanishes under $\Lambda_{\mathrm{dec}}$---are precisely those whose destruction generates entropy and reduces free energy.

We note that the relationship between coherence and extractable work is operationally subtle: the free energy $k_B T \cdot C_r$ is not in general directly extractable as work without additional resources such as a coherent reference frame \cite{lostaglio2015description,korzekwa2016extraction}.  The identity~\eqref{eq:free-energy-general} is exact as a free energy accounting; its operational interpretation requires care.

\begin{proposition}[Classical limit]
\label{prop:classical-limit}
When $\rho$ is already diagonal in the pointer basis $(\rho = \mathrm{diag}(p))$, $C_r(\rho) = 0$ and $\Delta F = 0$---no free energy is stored in coherences and none is destroyed by dephasing.
\end{proposition}

\begin{proof}
When $\rho$ is diagonal, $\Lambda_{\mathrm{dec}}[\rho] = \rho$, so $C_r(\rho) = S(\rho \,\|\, \rho) = 0$ and $F(\rho) = F(\Lambda_{\mathrm{dec}}[\rho])$.
\end{proof}

\begin{remark}[Classical coarse-graining]
\label{rem:classical-coarse-graining}
Classical coarse-graining (projecting a fine-grained probability distribution onto coarser bins) and quantum dephasing are complementary dimension-reducing processes: both map to lower-dimensional descriptions, but they operate in different directions (populations vs.\ coherences) and are not nested.  The classical Landauer bound \cite{landauer1961irreversibility,bennett1982thermodynamics} applies to the former; Proposition~\ref{prop:free-energy-coherence} applies to the latter.
\end{remark}

\begin{remark}[Cost scaling with system size]
\label{rem:cost-scaling}
For a $d$-dimensional system in a maximally coherent pure state with $H$ diagonal in the dephasing basis, the free energy destroyed by decoherence is $\Delta F = k_B T \ln d$.  Since the dimensional collapse eliminates $d^2 - d$ information-geometric dimensions (Theorem~\ref{thm:dim-collapse}), the cost per collapsed dimension is approximately
\begin{equation}
\label{eq:cost-per-dim}
\frac{\Delta F}{d^2 - d} \;=\; \frac{k_B T \ln d}{d^2 - d} \;\approx\; \frac{k_B T \ln d}{d^2},
\end{equation}
which decreases with system size.  The total free energy change $\Delta F = k_B T \ln d$ grows only logarithmically with Hilbert space dimension, while the number of destroyed dimensions grows quadratically.  This scaling is consistent with the observation that macroscopic decoherence is extremely fast \cite{zurek2003decoherence}: the per-dimension free energy change vanishes in the macroscopic limit.
\end{remark}

Figure~\ref{fig:dlb} illustrates the coherent free energy $\Delta F / k_B T = C_r$ as a function of dephasing strength for three qubit states.

\begin{figure}[H]
\centering
\includegraphics[width=0.7\textwidth]{fig4_quantum_dlb.pdf}
\caption{\textbf{Coherent free energy under progressive dephasing.}  $\Delta F / k_B T = C_r(\rho(\gamma))$ for a qubit with $H$ diagonal in the dephasing basis, as a function of the dephasing parameter $\gamma \in [0, 1/2]$.  A maximally coherent pure state $|{+}\rangle$ starts at $C_r = \ln 2$ and decays to zero at full dephasing; a partially coherent mixed state shows reduced but qualitatively similar behavior; an already-diagonal state has $C_r = 0$ throughout.}
\label{fig:dlb}
\end{figure}

\subsection{Pointer states and the geometry of einselection}
\label{sec:pointer}

The preceding analysis assumed a given pointer basis.  The decoherence channel depends on this choice: different bases project $\mathcal{M}_Q$ onto different embeddings of $\Delta_{d-1}$, annihilating the same \emph{number} of dimensions but \emph{different} dimensions.  Zurek's einselection \cite{zurek2003decoherence,zurek2025decoherence} resolves the ambiguity: the pointer states $\{|s_k\rangle\}$ are determined by the system--environment interaction Hamiltonian $H_{SE}$.  In our framework, einselection selects which $(d-1)$-dimensional submanifold the collapse targets; the collapse magnitude $d^2 - 1 \to d - 1$ is basis-independent (Corollary~\ref{cor:deff-collapse}).

Zurek's quantum Darwinism \cite{zurek2003decoherence,zurek2025decoherence,riedel2010zurek} adds that pointer states are \emph{replicated} across independent environmental fragments---Riedel and Zurek \cite{riedel2010zurek} showed $\sim\!10^7$ redundant copies of a dust grain's position are recorded in scattered photons within a microsecond.  In the dimensional collapse framework, each fragment acquires a copy of the same $(d-1)$-dimensional simplex.  This redundancy makes the classical projection robust: reversing it would require restoring phase relationships across all fragments simultaneously.

%% ============================================================
\section{Worked Examples}
\label{sec:examples}

We illustrate the framework with exactly solvable examples.

\subsection{Qubit dephasing}
\label{sec:qubit-dephasing}

For a single qubit ($d = 2$) parameterized by polar coordinates $(\theta, \phi)$ on the Bloch ball, partial dephasing with parameter $\gamma \in [0, 1/2]$ shrinks the equatorial components while preserving the polar component:
\begin{equation}
\label{eq:bloch-dephasing}
\mathbf{r} \;\mapsto\; \bigl((1-2\gamma)\sin\theta\cos\phi,\; (1-2\gamma)\sin\theta\sin\phi,\; \cos\theta\bigr).
\end{equation}
At $\gamma = 1/2$ the Bloch sphere collapses to the $z$-axis segment $\Delta_1$, as predicted by Theorem~\ref{thm:dim-collapse}.  The QFIM is $F_Q(\gamma) = \mathrm{diag}(1,\; (1-2\gamma)^2 \sin^2\theta)$ \cite{lu2013fisher}: the population parameter retains unit Fisher information while the phase parameter decays quadratically.  For equatorial states ($\theta = \pi/2$):
\begin{equation}
\label{eq:deff-qubit}
D_{\mathrm{eff}}^Q(\gamma) = \frac{\bigl(1 + (1-2\gamma)^2\bigr)^2}{1 + (1-2\gamma)^4},
\end{equation}
decreasing monotonically from $2$ to $1$ (Figure~\ref{fig:bloch}).

\begin{figure}[H]
\centering
\includegraphics[width=\textwidth]{fig1_bloch_collapse.pdf}
\caption{\textbf{Dimensional collapse of a qubit under dephasing.} Left: the Bloch sphere contracts to an oblate ellipsoid as $\gamma$ increases, collapsing to the $z$-axis at $\gamma = 1/2$. Right: $D_{\mathrm{eff}}^Q$ decreases from 2 (full quantum) to 1 (classical). Marked points correspond to the dephasing levels shown on the left.}
\label{fig:bloch}
\end{figure}

\subsection{\texorpdfstring{$N$}{N}-qubit register}
\label{sec:n-qubit}

For $N$ qubits ($d = 2^N$), the collapse ratio $1/(d+1) = 1/(2^N+1)$ vanishes exponentially (Figure~\ref{fig:fraction}): by $N = 10$ qubits, less than $0.1\%$ of the information-geometric degrees of freedom survive.

\begin{figure}[H]
\centering
\includegraphics[width=0.7\textwidth]{fig2_classical_fraction.pdf}
\caption{\textbf{Classical dimensionality ratio} $(d-1)/(d^2-1) = 1/(d+1)$ vs.\ number of qubits $N$ ($d = 2^N$).}
\label{fig:fraction}
\end{figure}

\subsection{Spin-boson model}
\label{sec:spin-boson}

We now consider continuous-time decoherence: a qubit coupled to an Ohmic bosonic bath \cite{leggett1987dynamics} with spectral density $J(\omega) = \eta\,\omega\, e^{-\omega/\omega_c}$ ($\eta$ = coupling strength, $\omega_c$ = cutoff).  In the high-temperature Markovian regime ($k_BT \gg \hbar\omega_c$), coherences decay as $\rho_{01}(t) = \rho_{01}(0)\, e^{-\Gamma(t)}$ with $\Gamma(t) = \eta k_BT\, t$.  Mapping to the qubit dephasing of Section~\ref{sec:qubit-dephasing} via $(1 - 2\gamma) \to e^{-\Gamma(t)}$:
\begin{equation}
\label{eq:deff-spinboson}
D_{\mathrm{eff}}^Q(t) = \frac{\bigl(1 + e^{-2\Gamma(t)}\bigr)^2}{1 + e^{-4\Gamma(t)}}, \qquad \tau_{\mathrm{dec}} \sim \frac{1}{2\eta k_BT}.
\end{equation}
After one decoherence time $\tau_{\mathrm{dec}}$, $D_{\mathrm{eff}}^Q$ has dropped roughly halfway from~$2$ to~$1$.  Figure~\ref{fig:trajectories} shows trajectories for several coupling strengths.

\begin{figure}[H]
\centering
\includegraphics[width=0.7\textwidth]{fig3_deff_trajectories.pdf}
\caption{\textbf{Dimensional collapse under spin-boson decoherence.} The effective quantum dimensionality $D_{\mathrm{eff}}^Q(t)$ for a qubit coupled to an Ohmic bosonic bath at high temperature, for four coupling strengths $\eta k_BT$. Weak coupling (0.01) preserves near-quantum dimensionality over the displayed time window; strong coupling (2.0) drives collapse to the classical value $D_{\mathrm{eff}}^Q = 1$ within $t \sim 1$. All curves share the same functional form but with different characteristic timescales.}
\label{fig:trajectories}
\end{figure}

At short times ($t \ll 1/\omega_c$), the exact decoherence function is quadratic, $\Gamma(t) \sim (\eta k_BT\, \omega_c)\, t^2$ \cite{leggett1987dynamics}, so the initial decay of $D_{\mathrm{eff}}^Q$ is quadratic rather than exponential---a signature of the quantum Zeno regime.  The crossover to Markovian (exponential) collapse occurs at $t \sim 1/\omega_c$.

%% ============================================================
\section{Discussion}
\label{sec:discussion}

\subsection{What this resolves}
\label{sec:resolves}

The quantum-to-classical transition has a precise geometric characterization: it is a dimensional collapse from $d^2 - 1$ to $d - 1$ in the quantum Fisher information manifold. The collapsed dimensions are specifically the quantum coherences---the $d^2 - d$ off-diagonal parameters that encode phase relationships between pointer states. The surviving dimensions are the $d - 1$ classical populations. When the Hamiltonian is diagonal in the dephasing basis, the free energy change under dephasing equals $k_B T \cdot C_r$ (Corollary~\ref{cor:diagonal-hamiltonian}).

The Bures separation angle (Theorem~\ref{thm:bures-angle}) adds a dynamical dimension to this picture: the metric separation between classical and quantum degrees of freedom is state-dependent and sharpened by dephasing.  The classical--quantum tangent split is oblique exactly when the state simultaneously carries nonzero coherence and nonzero population imbalance; dephasing monotonically restores orthogonality.  The clean geometric distinction of Theorem~\ref{thm:orthogonal-decomposition} is therefore realized progressively as decoherence proceeds, reaching perfect orthogonality at full dephasing.  The separation angle also carries estimation-theoretic meaning: it determines the fraction of SLD Fisher information for the population parameter that survives Schur-complement elimination of the coherence nuisance parameters (Corollary~\ref{cor:estimation}), connecting the geometric picture to multiparameter quantum estimation \cite{zhang2025hybrid,albarelli2019restoring,wang2024multiparameter}.

The qudit generalization reveals that the orthogonality criterion tightens sharply with Hilbert space dimension.  For qubits, equal populations suffice for orthogonality even in the presence of coherence.  For $d \geq 3$, Theorem~\ref{thm:spectator-orthogonality} proves that Bures-orthogonality fails whenever a nonzero coherence has a decoupled spectator index, covering all states with sparse coherence structure.  We conjecture that only diagonal states achieve orthogonality for $d \geq 3$ (Conjecture~\ref{conj:global-orthogonality}); a numerical survey of random states (Figure~\ref{fig:qudit-angles}) supports this for $d = 3, 4, 5$.

This characterization is substrate-independent: the collapse ratio $1/(d+1)$ depends only on the Hilbert space dimension~$d$, not on the system--environment coupling.  Physical parameters determine the \emph{rate} of collapse (Section~\ref{sec:spin-boson}); the \emph{magnitude} is fixed by the algebraic structure of quantum state space.

\subsection{What this does not resolve}
\label{sec:not-resolves}

This framework characterizes the geometric mechanism by which decoherence maps quantum state space onto classical probability space.  It does not address several deeper questions.

\emph{Why definite outcomes exist.}  The dimensional collapse framework shows \emph{how} a $(d^2-1)$-dimensional quantum manifold projects onto a $(d-1)$-dimensional classical simplex, but takes as given that quantum state spaces are high-dimensional and that observers access low-dimensional projections.  Why nature has this dimensional architecture---why the quantum--classical gap exists at all---is a question about the structure of the universe that lies outside empirical science.

\emph{The Born rule.} The QFIM and von Neumann entropy both presuppose the Born rule.  A derivation from more primitive principles \cite{zurek2003decoherence} would strengthen this framework but lies outside our scope.

\emph{The specific outcome.} Our framework characterizes the space of possible classical outcomes ($\Delta_{d-1}$) and the free energy change associated with reaching it, but does not explain why a particular outcome obtains.

\emph{The timing of commitment.} The dimensional collapse is continuous in $\gamma$ (or $t$); there is no critical threshold at which classicality abruptly appears.  Whether there is a ``moment of fact'' may require interpretive commitments beyond those made here.

\emph{The fully-coherent higher-dimensional case.} For $d \geq 3$, the spectator theorem (Theorem~\ref{thm:spectator-orthogonality}) and perturbative analysis (Theorem~\ref{thm:qudit-angles}) together cover all states except those with fully dense coherence far from the simplex.  Whether the orthogonality criterion holds globally (Conjecture~\ref{conj:global-orthogonality}) remains open; an adversarial optimization search finds no counterexample (Appendix~\ref{app:numerics}).

\subsection{Connection to the epiontic interpretation}
\label{sec:epiontic}

Zurek \cite{zurek2025decoherence} terms the quantum state ``epiontic''---simultaneously epistemic and ontic.  The dimensional collapse framework gives this a geometric reading: $d^2 - 1$ parameters are needed to specify the state (ontic), but after decoherence only $d - 1$ are accessible (epistemic).  The transition is not a change of philosophical stance but a geometric projection: the accessible information-geometric structure collapses from $d^2 - 1$ to $d - 1$ dimensions while unitarity is preserved at the level of system-plus-environment.

\subsection{Connection to dimensional projection theory}
\label{sec:projection-theory}

This work generalizes prior results on measurement as dimensional collapse in classical systems \cite{todd2025falsifiability,todd2025timing}: the classical Projection Bound (coarse-graining of populations) and the quantum free energy decomposition (dephasing of coherences) are complementary bounds that are not nested (Remark~\ref{rem:classical-coarse-graining}).

Two directions for future work emerge naturally.  First, quantum error correction resists dimensional collapse by encoding logical qubits in subspaces insensitive to dominant error operators; the present framework may provide a geometric characterization of the error-correction threshold.  Second, extending beyond complete dephasing to amplitude damping and general CPTP maps would characterize collapse to intermediate manifolds where $d - 1 < D_{\mathrm{eff}}^Q < d^2 - 1$.

The principal open problem is Conjecture~\ref{conj:global-orthogonality}: proving that for $d \geq 3$, the fully-coherent case (complete coherence graph) also breaks Bures-orthogonality.  The spectator theorem (Theorem~\ref{thm:spectator-orthogonality}) and perturbative analysis (Theorem~\ref{thm:qudit-angles}) together cover all states except those with fully dense coherence far from the simplex; closing this gap would complete the higher-dimensional theory.

\bibliographystyle{unsrt}
\bibliography{references}

%% ============================================================
\appendix
\section{Numerical Methods}
\label{app:numerics}

The numerical results in Figure~\ref{fig:qudit-angles} and the conjecture-support discussion use the following procedures.

\paragraph{Random state generation.}
Full-rank density matrices are drawn from the Hilbert--Schmidt measure: generate a $d \times d$ matrix $A$ with i.i.d.\ standard complex Gaussian entries, form $\rho = A A^\dagger / \mathrm{Tr}[A A^\dagger]$.  This distribution is unitarily invariant and concentrates on full-rank states.

\paragraph{SLD computation.}
For each tangent direction $X$ (traceless Hermitian), the symmetric logarithmic derivative $L_X$ satisfying $\rho L_X + L_X \rho = 2X$ is computed in the eigenbasis of~$\rho$: $(L_X)_{mn}^{\mathrm{eig}} = 2 X_{mn}^{\mathrm{eig}} / (\lambda_m + \lambda_n)$, then transformed back to the pointer basis.  Full-rank states ensure $\lambda_m + \lambda_n > 0$ for all~$m, n$.

\paragraph{Principal angles.}
The $d - 1$ principal angles between $T_\rho^{\mathrm{diag}}$ (dimension $d - 1$) and $T_\rho^{\mathrm{off}}$ (dimension $d^2 - d$) are computed as $\theta_a = \arccos(\sigma_a)$, where $\sigma_1, \ldots, \sigma_{d-1}$ are the singular values of $G_{DD}^{-1/2}\, G_{DO}\, G_{OO}^{-1/2}$.  Here $G_{DD}$, $G_{OO}$, and $G_{DO}$ are the Gram and cross-Gram matrices of the respective bases under the Bures metric.

\paragraph{Adversarial search.}
To test whether any non-diagonal state can approach orthogonality, we perform a local optimization for each~$d$.  Full-rank density matrices are parameterized by $d + d(d-1)$ real parameters: softmax log-weights for the diagonal and scaled real/imaginary parts for the off-diagonal entries.  We run $100$ independent L-BFGS-B optimizations (each initialized from a random state) that maximize the minimum principal angle $\min_a \theta_a$.  The best achieved angles are reported in Figure~\ref{fig:qudit-angles} as vertical lines; in no case does the optimum exceed the maximum found by random sampling by more than a few degrees, and no optimized state approaches $90^\circ$.

\paragraph{Software.}
All computations use NumPy (linear algebra) and SciPy (optimization).  The complete code is available in the repository at \texttt{code/generate\_figures.py}.

\end{document}
